% ── preliminaries.tex · §2 Preliminaries (PP & LP) ─────────────────
\section{Preliminaries}\label{sec:prelim}
\begin{figure}[t]
	\centering
	\begin{minipage}[c]{0.40\linewidth}
		\begin{codePP}
[ALL2] <not(!x.!y.not(x = x))>
[ALL6] <not(!(x,y).not(x = x))>
[STOP_1] <FAUX>
[EVR2_1] <not(x = x)>
[ALL7(1)] <!(x,y).not(x = x) => FAUX>
[FIN(FAUX)] <$\ldots$>
[NRM1] <forall(x,y).FAUX => FAUX>
[STOP_NORM] <FAUX => FAUX>
[NRM] <forall(x,y).FAUX => FAUX>
[FIN(FAUX => FAUX)] <$\ldots$>
[AXM7] <FAUX => FAUX>
		\end{codePP}
	\end{minipage}\hfill
	\begin{minipage}[c]{0.57\linewidth}
		\centering\footnotesize
		\begin{prooftree}
			\infer0[\rn{STOP\_1}]{(\seq \bot) \yields \bot}
			\infer1[\rn{EVR2\_1}]{(\seq \lnot(x = x)) \yields \bot}
			\infer0[\rn{AXM7}]{\seq \bot \Rightarrow \bot}
			\infer1[\rn{NRM1}]{\seq (\Dia (x,y) \cdot \bot) \Rightarrow \bot}
			\infer2[\rn{ALL7}]{\seq (\forall (x,y) \cdot \lnot(x = x)) \Rightarrow \bot}
			\infer1[\rn{ALL6}]{\seq \lnot \forall (x,y) \cdot \lnot(x = x)}
			\infer1[\rn{ALL2}]{\seq \lnot \forall x \cdot \forall y \cdot \lnot(x = x)}
		\end{prooftree}
	\end{minipage}
	\caption{Replay of benchmark \texttt{og/22} and its proof tree.}
	\label{fig:replay}
\end{figure}

\paragraph{The Predicate Prover.}\label{sec:pp-syntax}
%
The syntax of \pp{} is given by the following grammar:
\begin{align*}
	E, F & \Coloneqq x \mid n \mid b \mid E \mapsto F \mid F(E_1,\ldots,E_k)
	\mid E + F \mid E - F \mid -E                                     \\
	P, Q & \Coloneqq E \mid E = F \mid E \leq F
	\mid E_1,\ldots,E_k \in F
	\mid \lnot P \mid P \land Q \mid P \lor Q \mid P \Rightarrow Q
	\mid P \Leftrightarrow Q
	\mid \mathcal{Q}\,\vec{x} \cdot P
\end{align*}
where $x$ ranges over identifiers, and $n$ over numerals, and $b$ over booleans.
%
\pp's \emph{expressions} $E, F$ are variables, numerals, pairs $E \mapsto F$,
function applications $F(E_1,\ldots,E_k)$, and linear arithmetic ($+$, $-$);
its \emph{predicates} $P, Q$ are membership $E_1,\ldots,E_k \in F$, equality,
$\leq$, coerced (boolean) expressions $E$, propositional connectives,
and quantification $\mathcal{Q}\,\vec{x} \cdot P$ with
$\mathcal{Q} \in \{\forall, \exists, \Dia, \Hrt\}$.
%
Quantifiers bind a list of variables; $\forall (x,y) \cdot P$ is distinct from
$\forall x \cdot \forall y \cdot P$, and dedicated rules regroup the latter
within proofs.
%
The quantifiers $\Dia$ and $\Hrt$ are idiosyncratic to PP; semantically both
are $\forall$ and are used only to guide PP's proof search.

\pp{} proves \emph{sequents} $H \seq P$, where $H$ is a list of
hypothesis predicates and $P$ a goal predicate.
%
The inference rules for deriving sequents are specified in the annex of its
internal documentation~\cite{pp_spec}.
%
There are roughly \numRulesSpec{} rules, grouped into families:
propositional (\rn{AND}, \rn{OR}, \rn{IMP}, \rn{EQV},
\rn{NOT}), hypothesis discharge (\rn{AXM1--9}), quantifiers
(\rn{ALL1--9}, \rn{XST1--8}), equality
(\rn{EVR}, \rn{EAXM}, \rn{EQC}, \rn{EQS}, \rn{ECTR}, \ldots), booleans
(\rn{BOOL}), linear arithmetic (\rn{AR1--13}),
and normalisation (\rn{NRM1--30}), among others.
%
Each rule reduces the current goal to zero, one, or two
premises, possibly under computational side-conditions
(e.g., ``$P$ occurs in the list of hypotheses $H$'',
``$a$ and $b$ are numerals such that $a > b$'').

Alongside sequents, proofs manipulate a second judgment form known as a
\emph{result}; essentially an equality proof used by \pp{}'s
normalization procedure.
%
We write $(H \seq P) \yields [Q_1, \ldots, Q_n]$ for
``under $H$, $P$ normalises to the list $[Q_1, \ldots, Q_n]$''.
The conjunction $Q_1 ∧ \ldots ∧ Q_n$ is the normal form of $P$.
%
With a few exceptions, each the result form of each inference rule
(written with a $\mathsf{\_1}$ suffix) is derived by one of three schemas:
the empty list (the formula is discharged),
a singleton $[Q]$ (a single transformed formula),
or the concatenation of its two premises' lists (a combination).
%
For example, consider the \rn{AND4} rule and its result form \rn{AND4\_1}:
\[
\begin{prooftree}
\hypo{H \seq Q}
\hypo{H \seq P}
\infer2[\rn{AND4}]{H \seq P \land Q}
\end{prooftree}
\qquad
\begin{prooftree}
\hypo{(H \seq Q) \yields R_1}
\hypo{(H \seq P) \yields R_2}
\infer2[\rn{AND4\_1}]{(H \seq P \land Q) \yields R_1 \app R_2}
\end{prooftree}
\]
%
The only rules that use a result derivation to derive a sequent
are \rn{ALL7} and \rn{XST8}, given below:
\[
\begin{array}{c}
\begin{prooftree}
\hypo{(H \seq P) \yields [R_1, \ldots, R_n]}
\hypo{H \seq (\Dia x \cdot R_1 \land \ldots \land R_n) \Rightarrow G}
\infer2[\rn{ALL7}]{H \seq (\forall x \cdot P) \Rightarrow G}
\end{prooftree}
\\[5mm]
\begin{prooftree}
\hypo{(H \seq \lnot P) \yields [R_1, \ldots, R_n]}
\hypo{H \seq (\Dia x \cdot R_1 \land \ldots \land R_n) \Rightarrow \bot}
\infer2[\rn{XST8}]{H \seq \exists x \cdot P}
\end{prooftree}
\end{array}
\]
%
\noindent
\autoref{fig:replay} shows an example of the \replay{} format.
%
Each replay line carries a rule name and, in parentheses, an optional
rule-specific argument.
%
A rule's line precedes the lines of its subproofs, with two exceptions.
%
First, the result derivation feeding \rn{ALL7} or \rn{XST8} appears
before the rule that consumes it.
%
Second, the replay interleaves marker lines that do not correspond to
inference rules: \texttt{NRM} and \texttt{STOP\_NORM}, which bookend
\pp's normalisation passes, and \texttt{FIN}, which carries the predicate
such a pass produced.
%

\paragraph{\lp}\label{sec:lp}
%
\lp~\cite{hondet:hal-02981561} is a proof assistant for the
\lpcmr~\cite{CousineauDowek}.
%
%
Unlike \dedukti{}~\cite{assaf2023dedukti}, a bare type-checker for the same
calculus, \lp{} has implicit arguments and an interactive tactic language.
%
The term syntax is given as follows:
\[
	t, u \Coloneqq \TYPE \mid \KIND \mid x \mid c
	\mid t\,u \mid \lambda\,x : t,\, u \mid \Pi\,x : t,\, u
\]
where $x$ ranges over variables, $c$ ranges over declared constants,
and terms are identified up to $\alpha$-equivalence.
%
Typing is standard for the \lpcmr{} (the full system is in
\autoref{app:typing}); its \emph{definitional equality} (i.e., $A \Equiv B$)
is generated by $\beta$-equivalence and user-declared rewrite rules.
%
A \lp{} development is a sequence of declarations:
\[
  \begin{array}{rcl}
    d &\Coloneqq& m^{*}\;\text{\lpkw{symbol}}\;x\;p^{*} : t\;[\,≔ b\,]
        \;\mid\; \text{\lpkw{rule}}\;\ell ↪ r
        \;\mid\; \text{\lpkw{inductive}}\;x\;p^{*} : t ≔ \mathit{cons} \\
    p &\Coloneqq& (x : t) \;\mid\; [x : t] \\
    m &\Coloneqq& \text{\lpmod{constant}} \;\mid\; \text{\lpmod{injective}} \;\mid\; \text{\lpmod{opaque}} \\
    b &\Coloneqq& t \;\mid\; \text{\lpkw{begin}}\;\mathit{tac}^{*}\;\text{\lpkw{end}}
  \end{array}
\]
The \lpkw{symbol} command introduces a typed constant
with explicit $(x : A)$ and implicit $[x : A]$ parameters $p$
and an optional definition $b$; \lpkw{rule} adds a
type-preserving rewrite rule over pattern variables \texttt{\$x};
\lpkw{inductive} generates an inductive type with an induction principle.
%
Modifiers $m$ refine a symbol: \lpmod{constant} symbols admit no rewrite rule at
their head, \lpmod{injective} ones may be inverted during unification, and an
\lpmod{opaque} symbol's definition is sealed after type-checking.
%
A definition $b$ is a term or a proof script (\lpkw{begin} \ldots{}
\lpkw{end}) whose tactics (e.g., \lpkw{assume}, \lpkw{refine},
\lpkw{rewrite}, \lpkw{have}, \lpkw{induction}) elaborate to a proof term.
%
The grammars above are abridged; \autoref{app:syntax} gives an extended grammar.
